\documentclass[landscape,final,archE1,fontscale=0.32]{baposter}

\usepackage{calc}
\usepackage{graphicx}
\usepackage{amsmath}
\usepackage{amssymb}
\usepackage{relsize}
\usepackage{rotating}
\usepackage{bm}
\usepackage{url}
\usepackage{setspace}

\usepackage{tabularx}
\usepackage{graphicx}

\usepackage{palatino}
\usepackage{booktabs, multicol, multirow}

\usepackage[margin=1em]{caption}
% package I added
\usepackage{sidecap}
\usepackage{floatrow}
\usepackage{wrapfig}


\bibliographystyle{plain}

\graphicspath{{images/}{../images/}}
\usetikzlibrary{calc}

\newcommand{\SET}[1]  {\ensuremath{\mathcal{#1}}}
\newcommand{\MAT}[1]  {\ensuremath{\boldsymbol{#1}}}
\newcommand{\VEC}[1]  {\ensuremath{\boldsymbol{#1}}}
\newcommand{\Video}{\SET{V}}
\newcommand{\video}{\VEC{f}}
\newcommand{\track}{x}
\newcommand{\Track}{\SET T}
\newcommand{\LMs}{\SET L}
\newcommand{\lm}{l}
\newcommand{\PosE}{\SET P}
\newcommand{\posE}{\VEC p}
\newcommand{\negE}{\VEC n}
\newcommand{\NegE}{\SET N}
\newcommand{\Occluded}{\SET O}
\newcommand{\occluded}{o}

%%%%%%%%%%%%%%%%%%%%%%%%%%%%%%%%%%%%%%%%%%%%%%%%%%%%%%%%%%%%%%%%%%%%%%%%%%%%%%%%
%%%% Some math symbols used in the text
%%%%%%%%%%%%%%%%%%%%%%%%%%%%%%%%%%%%%%%%%%%%%%%%%%%%%%%%%%%%%%%%%%%%%%%%%%%%%%%%

%%%%%%%%%%%%%%%%%%%%%%%%%%%%%%%%%%%%%%%%%%%%%%%%%%%%%%%%%%%%%%%%%%%%%%%%%%%%%%%%
% Multicol Settings
%%%%%%%%%%%%%%%%%%%%%%%%%%%%%%%%%%%%%%%%%%%%%%%%%%%%%%%%%%%%%%%%%%%%%%%%%%%%%%%%
\setlength{\columnsep}{1.5em}
\setlength{\columnseprule}{0mm}

%%%%%%%%%%%%%%%%%%%%%%%%%%%%%%%%%%%%%%%%%%%%%%%%%%%%%%%%%%%%%%%%%%%%%%%%%%%%%%%%
% Save space in lists. Use this after the opening of the list
%%%%%%%%%%%%%%%%%%%%%%%%%%%%%%%%%%%%%%%%%%%%%%%%%%%%%%%%%%%%%%%%%%%%%%%%%%%%%%%%
\newcommand{\compresslist}{%
\setlength{\itemsep}{1pt}%
\setlength{\parskip}{0pt}%
\setlength{\parsep}{0pt}%
}

%%%%%%%%%%%%%%%%%%%%%%%%%%%%%%%%%%%%%%%%%%%%%%%%%%%%%%%%%%%%%%%%%%%%%%%%%%%%%%
%%% Begin of Document
%%%%%%%%%%%%%%%%%%%%%%%%%%%%%%%%%%%%%%%%%%%%%%%%%%%%%%%%%%%%%%%%%%%%%%%%%%%%%%

\begin{document}

%%%%%%%%%%%%%%%%%%%%%%%%%%%%%%%%%%%%%%%%%%%%%%%%%%%%%%%%%%%%%%%%%%%%%%%%%%%%%%
%%% Here starts the poster
%%%---------------------------------------------------------------------------
%%% Format it to your taste with the options
%%%%%%%%%%%%%%%%%%%%%%%%%%%%%%%%%%%%%%%%%%%%%%%%%%%%%%%%%%%%%%%%%%%%%%%%%%%%%%
% Define some colors

\definecolor{lightblue}{rgb}{0.145,0.6666,1}

\hyphenation{resolution occlusions}
%%
\begin{poster}%
  % Poster Options
  {
  % Show grid to help with alignment
  grid=false,
  columns=3,
  % Column spacing
  colspacing=1.6em,
  % Color style
  bgColorOne=white,
  bgColorTwo=white,
  borderColor=lightblue,
  headerColorOne=black,
  headerColorTwo=lightblue,
  headerFontColor=white,
  boxColorOne=white,
  boxColorTwo=lightblue,
  % Format of textbox
  textborder=rectangle,
  % Format of text header
  eyecatcher=true,
  headerborder=closed,
  headerheight=0.12\textheight,
%  textfont=\sc, An example of changing the text font
  headershape=roundedright,
  headershade=shadelr,
  headerfont=\Large\bf, %Sans Serif
  textfont={\setlength{\parindent}{1.5em}},
  boxshade=plain,
%  background=shade-tb,
  background=plain,
  linewidth=2pt
  }
  % ICSI logo
  {\includegraphics[height=7.5em]{images/ICSI_color}}
  % Title
  {\bf{Parallelization of Neural Network Training using \\Model Averaging and Butterfly Mixing}}
  % Authors
  {Hang Su$^1$$^,$$^2$, Haoyu Chen$^2$\\
  $^1$ International Computer Science Institute, CA, USA ~~~ $^2$ Dept. of EECS, UC Berkeley, CA, USA }
  % University logo
  {% The makebox allows the title to flow into the logo, this is a hack because of the L shaped logo.
    \includegraphics[height=7.5em]{images/ucbseal}
  }

%%%%%%%%%%%%%%%%%%%%%%%%%%%%%%%%%%%%%%%%%%%%%%%%%%%%%%%%%%%%%%%%%%%%%%%%%%%%%%
%%% Now define the boxes that make up the poster
%%%---------------------------------------------------------------------------
%%% Each box has a name and can be placed absolutely or relatively.
%%% The only inconvenience is that you can only specify a relative position
%%% towards an already declared box. So if you have a box attached to the
%%% bottom, one to the top and a third one which should be in between, you
%%% have to specify the top and bottom boxes before you specify the middle
%%% box.
%%%%%%%%%%%%%%%%%%%%%%%%%%%%%%%%%%%%%%%%%%%%%%%%%%%%%%%%%%%%%%%%%%%%%%%%%%%%%%
    %
    % A coloured circle useful as a bullet with an adjustably strong filling
    \newcommand{\colouredcircle}{%
      \tikz{\useasboundingbox (-0.2em,-0.32em) rectangle(0.2em,0.32em); \draw[draw=black,fill=lightblue,line width=0.03em] (0,0) circle(0.18em);}}

%%%%%%%%%%%%%%%%%%%%%%%%%%%%%%%%%%%%%%%%%%%%%%%%%%%%%%%%%%%%%%%%%%%%%%%%%%%%%%
\headerbox{Goal}{name=goal,column=0,row=0}{
%%%%%%%%%%%%%%%%%%%%%%%%%%%%%%%%%%%%%%%%%%%%%%%%%%%%%%%%%%%%%%%%%%%%%%%%%%%%%%
\begin{itemize}
\setlength \itemsep{0.2em}
  \item Distributed neural network training using multiple GPUs.
  \item Model average for distributed network training.
  \item Apply different network communication strategies (all-reduce, butterfly and ring).
\end{itemize}
}

%%%%%%%%%%%%%%%%%%%%%%%%%%%%%%%%%%%%%%%%%%%%%%%%%%%%%%%%%%%%%%%%%%%%%%%%%%%%%%
\headerbox{Background}{name=background,column=0,below=goal}{
%%%%%%%%%%%%%%%%%%%%%%%%%%%%%%%%%%%%%%%%%%%%%%%%%%%%%%%%%%%%%%%%%%%%%%%%%%%%%%
\subsection*{Condition}
\begin{itemize}
\setlength \itemsep{0.2em}
\item Large scale continuous time automatic speech recognition task.
\item Deep Neural Network with fully connected layers and thousands of hidden nodes.
\item Multiple GPUs available on different computing nodes.
\item Data communication among the nodes are bottleneck for distributed training.
\end{itemize}

\subsection*{Approaches}
\begin{itemize}
\setlength \itemsep{0.2em}
\item Data parallelism -- train neural networks locally using partitioned data (GPU based).
\item Average model parameters among computing nodes every few minibatches (OpenMPI based).
\item Apply butterfly / ring mixing network communication strategy to reduce data communication.
\item Special training algorithm (natural gradient SGD) for better covergence rate.
\end{itemize}
}

%%%%%%%%%%%%%%%%%%%%%%%%%%%%%%%%%%%%%%%%%%%%%%%%%%%%%%%%%%%%%%%%%%%%%%%%%%%%%%
\headerbox{Setup}{name=setup,column=0,above=bottom,below=background}{
%%%%%%%%%%%%%%%%%%%%%%%%%%%%%%%%%%%%%%%%%%%%%%%%%%%%%%%%%%%%%%%%%%%%%%%%%%%%%%
\begin{itemize}
\setlength \itemsep{0.2em}
\item OpenMPI for data communicate between nodes.
\item GPU for local neural network training (CUDA M1060 * 16)
\end{itemize}
\begin{center}
\includegraphics[width=0.45\linewidth]{framework.jpg}
\end{center}
}
%%%%%%%%%%%%%%%%%%%%%
%\headerbox{}{name=setup2, column = 1, row = 0}{

%The figure for the system's framework. 
%\begin{wrapfigure}{r}{0.5\textwidth}
%  \vspace{-30pt}
%  \begin{center}
%    \includegraphics[width=1.0\textwidth]{framework.jpg}
%  \end{center}
%  \vspace{-20pt}
  
%  \vspace{-10pt}
%\end{wrapfigure}
%When we plan to communicate the parameters, we need to pass the message from GPU to CPU, and send message from one CPU to another CPU by OpenMPI. 

%}
%%%%%%%%%%%%%%%%%%%%%%%%%%%%%%%%%%%%%%%%%%%%%%%%%%%%%%%%%%%%%%%%%%%%%%%%%%%%%%
\headerbox{Parallelization Strategy}{name=strategy,column=1,row=0}{
%%%%%%%%%%%%%%%%%%%%%%%%%%%%%%%%%%%%%%%%%%%%%%%%%%%%%%%%%%%%%%%%%%%%%%%%%%%%%%
\begin{multicols}{2}
\subsection*{All-reduce network}
After every $n$ local updates, all the nodes reduce their parameter together (average), 
and broadcast the parameter back to all the nodes. This requires high communication
bandwidth.

\vspace{-0.2cm}

\subsection*{Butterfly-mixing}
In Butterfly-mixing strategy, one node only communicates with one other node, averaging the 
recieved model with its own model. After $log(nodeSize)$ communication, the information 
is fully communicated.

\vspace{-0.2cm}

\subsection*{Ring network}
In ring network, each node sends its parameter to its right neighbor and receive parameters 
from its left neighbor. After $nodeSize$ communication, the information is fully communicated.
This method benefits from the locality of computing nodes.

\begin{center}
  \includegraphics[width=0.40\textwidth]{allreduce.jpg}
\end{center}

\begin{center}
  \includegraphics[width=0.4\textwidth]{butterfly.jpg}
\end{center}

\begin{center}
  \includegraphics[width=0.4\textwidth]{ring.jpg}
\end{center}
\end{multicols}

}

\headerbox{Natural Gradient SGD}{name=sgd,column=1,below=strategy, above=bottom}{
\begin{itemize}
\setlength \itemsep{0.2em}
\item Take steps along the gradient of a Riemannian parameter surface.
\item Use inverse fisher matrix $E_t$ as learning-rate matrix.
\item $\alpha$ denotes original learning rate, $g_t$ denotes plain gradient. Update
parameters using:
\end{itemize}
\begin{equation}
\theta_{t+1} = \theta_{t} + \alpha_t E_t g_t
\end{equation}
}

%%%%%%%%%%%%%%%%%%%%%%%%%%%%%%%%%%%%%%%%%%%%%%%%%%%%%%%%%%%%%%%%%%%%%%%%%%%%%%
\headerbox{Experiments}{name=experiments,column=2,row=0}{
%%%%%%%%%%%%%%%%%%%%%%%%%%%%%%%%%%%%%%%%%%%%%%%%%%%%%%%%%%%%%%%%%%%%%%%%%%%%%%
\subsection*{Convergence Status}
% a graph of objective function vs iterations
% natural gradient vs normal gradient 
\begin{center}
  \captionof{table}{Convergence status on CV set}
  \begin{tabular}{c|c|c|c|c|c|c}
    \hline
           \multicolumn{2}{c|}{Nodes}         & 1    & 2    & 4    & 8    & 16 \\
    \hline
\multirow{2}{*}{Allreduce} &    Obj (loss)    & 1.43 & 1.57 & 1.45 & 1.45 & 1.48\\
                           &    Frame Acc(\%) & 59.7 & 59.3 & 60.2 & 59.9 & 58.5\\
    \hline
\multirow{2}{*}{Butterfly} &    Obj (loss)    & --   & --   & 1.47 & 1.52 & 1.70\\
                           &    Frame Acc(\%) & --   & --   & 60.0 & 58.4 & 54.7 \\
    \hline
\multirow{2}{*}{Ring}      &    Obj (loss)    & --   & --   & 1.47 & 1.62 & 1.96\\
                           &    Frame Acc(\%) & --   & --   & 59.8 & 56.2 & 50.7 \\
    \hline
  \end{tabular}
\end{center}


\subsection*{Scaling factor}
\begin{center}
  \captionof{table}{Speedup and scaling factor}
  \begin{tabular}{c|c|c|c|c|c|c}
    \hline
           \multicolumn{2}{c|}{Nodes}   & 1    & 2    & 4    & 8    & 16 \\
    \hline
\multirow{2}{*}{Allreduce} &    Time    & 6.59 & 3.39 & 1.90 & 1.01 & 0.700\\
                           &    speedup & --   & 1.94 & 3.47 & 6.52 & 9.41 \\
    \hline
\multirow{2}{*}{Butterfly} &    Time    & --   & --   & 1.83 & 1.00 & 0.582\\
                           &    speedup & --   & --   & 3.60 & 6.59 & 11.32\\
    \hline
\multirow{2}{*}{Ring}      &    Time    & --   & --   & 1.84 & 1.01 & 0.517\\
                           &    speedup & --   & --   & 3.58 & 6.52 & 12.75\\
    \hline
  \end{tabular}
\end{center}

\subsection*{Recognition Results}
\begin{center}
  \captionof{table}{Word error rate on different setup}
  \begin{tabular}{c|c|c|c|c|c}
    \hline
    Nodes       & 1    & 2     & 4    & 8    & 16 \\
    \hline
    allreduce   & 18.4 & 17.9  & 18.3 & 18.3 & 18.7 \\
    butterfly   & --   & --    & 18.5 & 18.7 & 20.6 \\
    ring        & --   & --    & 18.5 & 19.6 & 22.1\\
    \hline
  \end{tabular}
\end{center}

}

%%%%%%%%%%%%%%%%%%%%%%%%%%%%%%%%%%%%%%%%%%%%%%%%%%%%%%%%%%%%%%%%%%%%%%%%%%%%%%
\headerbox{Conclusion}{name=conclusion,column=2,below=experiments,above=bottom}{
%%%%%%%%%%%%%%%%%%%%%%%%%%%%%%%%%%%%%%%%%%%%%%%%%%%%%%%%%%%%%%%%%%%%%%%%%%%%%%
\begin{itemize}
\setlength \itemsep{0.2em}
\item Neural network training can be efficiently speed up by using model averaging techniques.
\item Natural gradient method contributes to the convergence rate in model averaging setup.
\item Butterfly-mixing and ring reduction can reduce communication time in a further step.
\item Training using Butterfly-mixing and ring reduction does not converge well as the number of jobs goes up to 16.
\end{itemize}
}

\end{poster}

\end{document}
